% GENERATED FILE. Edit canonical lessons, then run npm run generate:books.
% canonical-source-hash: fnv1a64:7ec369ada08b0928
% canonical-lessons: TE-C127-naluka, TE-C127-mokalu, TE-C127-bhujam, TE-C127-kanubomma, TE-C127-bugga

\chapter{Tongue, Knee, Shoulder, Eyebrow, Cheek}
\label{ch:te-tongue-knee-shoulder-eyebrow-cheek}
\hlchaptermodality{127}

\begin{chapteropening}
\textbf{By the end of this chapter:} \emph{I can say the tongue, a knee, a shoulder, an eyebrow and a cheek.}

Complete the last lesson of chapter 127: i can say the tongue, a knee, a shoulder, an eyebrow and a cheek.
\end{chapteropening}

\section[\texorpdfstring{\te{నాలుక} (nāluka)}{నాలుక (nāluka)}]{\te{నాలుక} (nāluka) — the tongue}

\label{lesson:TE-C127-naluka}



\begin{quote}
Before the new one: say the Telugu for a camel, then the Telugu for the face.
\end{quote}

\subsection*{You'll want to know: \te{నాలుక}}
\textbf{\te{నాలుక}} — \emph{nāluka} — \textquotedblleft{}the tongue\textquotedblright{}.

\begin{grammarlens}[title={What you've built}]
One more word for this chapter.
\end{grammarlens}

\subsection*{Your turn}
\begin{itemize}
\raggedright
  \item \emph{Say it:} \emph{nāluka}
  \item \emph{Say it:} \emph{nāluka}, once more
  \item \emph{Say it:} say \emph{nāluka}
  \item \emph{Recall:} say the Telugu for a camel, then the Telugu for the face, then say \emph{nāluka} again
\end{itemize}

\begin{tcolorbox}[breakable,colback=teal!4,colframe=teal!35!black,title={Before you move on}]
What does \te{నాలుక} mean? (\textquotedblleft{}The tongue\textquotedblright{}.) Say it once more.
\end{tcolorbox}
\section[\texorpdfstring{\te{మోకాలు} (mōkālu)}{మోకాలు (mōkālu)}]{\te{మోకాలు} (mōkālu) — a knee}

\label{lesson:TE-C127-mokalu}



\begin{quote}
Before the new one: say the Telugu for the face, then the Telugu for the tongue.
\end{quote}

\subsection*{You'll want to know: \te{మోకాలు}}
\textbf{\te{మోకాలు}} — \emph{mōkālu} — \textquotedblleft{}a knee\textquotedblright{}.

\begin{grammarlens}[title={What you've built}]
One more word for this chapter.
\end{grammarlens}

\subsection*{Your turn}
\begin{itemize}
\raggedright
  \item \emph{Say it:} \emph{mōkālu}
  \item \emph{Say it:} \emph{mōkālu}, once more
  \item \emph{Say it:} say \emph{mōkālu}
  \item \emph{Recall:} say the Telugu for the face, then the Telugu for the tongue, then say \emph{mōkālu} again
\end{itemize}

\begin{tcolorbox}[breakable,colback=teal!4,colframe=teal!35!black,title={Before you move on}]
What does \te{మోకాలు} mean? (\textquotedblleft{}A knee\textquotedblright{}.) Say it once more.
\end{tcolorbox}
\section[\texorpdfstring{\te{భుజం} (bhujaṁ)}{భుజం (bhujaṁ)}]{\te{భుజం} (bhujaṁ) — a shoulder}

\label{lesson:TE-C127-bhujam}



\begin{quote}
Before the new one: say the Telugu for the tongue, then the Telugu for a knee.
\end{quote}

\subsection*{You'll want to know: \te{భుజం}}
\textbf{\te{భుజం}} — \emph{bhujaṁ} — \textquotedblleft{}a shoulder\textquotedblright{}.

\begin{grammarlens}[title={What you've built}]
One more word for this chapter.
\end{grammarlens}

\subsection*{Your turn}
\begin{itemize}
\raggedright
  \item \emph{Say it:} \emph{bhujaṁ}
  \item \emph{Say it:} \emph{bhujaṁ}, once more
  \item \emph{Say it:} say \emph{bhujaṁ}
  \item \emph{Recall:} say the Telugu for the tongue, then the Telugu for a knee, then say \emph{bhujaṁ} again
\end{itemize}

\begin{tcolorbox}[breakable,colback=teal!4,colframe=teal!35!black,title={Before you move on}]
What does \te{భుజం} mean? (\textquotedblleft{}A shoulder\textquotedblright{}.) Say it once more.
\end{tcolorbox}
\section[\texorpdfstring{\te{కనుబొమ్మ} (kanubomma)}{కనుబొమ్మ (kanubomma)}]{\te{కనుబొమ్మ} (kanubomma) — an eyebrow}

\label{lesson:TE-C127-kanubomma}



\begin{quote}
Before the new one: say the Telugu for a knee, then the Telugu for a shoulder.
\end{quote}

\subsection*{You'll want to know: \te{కనుబొమ్మ}}
\textbf{\te{కనుబొమ్మ}} — \emph{kanubomma} — \textquotedblleft{}an eyebrow\textquotedblright{}.

\begin{grammarlens}[title={What you've built}]
One more word for this chapter.
\end{grammarlens}

\subsection*{Your turn}
\begin{itemize}
\raggedright
  \item \emph{Say it:} \emph{kanubomma}
  \item \emph{Say it:} \emph{kanubomma}, once more
  \item \emph{Say it:} say \emph{kanubomma}
  \item \emph{Recall:} say the Telugu for a knee, then the Telugu for a shoulder, then say \emph{kanubomma} again
\end{itemize}

\begin{tcolorbox}[breakable,colback=teal!4,colframe=teal!35!black,title={Before you move on}]
What does \te{కనుబొమ్మ} mean? (\textquotedblleft{}An eyebrow\textquotedblright{}.) Say it once more.
\end{tcolorbox}
\section[\texorpdfstring{\te{బుగ్గ} (bugga)}{బుగ్గ (bugga)}]{\te{బుగ్గ} (bugga) — a cheek}

\label{lesson:TE-C127-bugga}



\begin{quote}
Before the new one: say the Telugu for a shoulder, then the Telugu for an eyebrow.
\end{quote}

\subsection*{You'll want to know: \te{బుగ్గ}}
\textbf{\te{బుగ్గ}} — \emph{bugga} — \textquotedblleft{}a cheek\textquotedblright{}.

\begin{grammarlens}[title={What you've built}]
One more word for this chapter.
\end{grammarlens}

\subsection*{Your turn}
\begin{itemize}
\raggedright
  \item \emph{Say it:} \emph{bugga}
  \item \emph{Say it:} \emph{bugga}, once more
  \item \emph{Say it:} say \emph{bugga}
  \item \emph{Recall:} say the Telugu for a shoulder, then the Telugu for an eyebrow, then say \emph{bugga} again
\end{itemize}

\begin{tcolorbox}[breakable,colback=teal!4,colframe=teal!35!black,title={Before you move on}]
What does \te{బుగ్గ} mean? (\textquotedblleft{}A cheek\textquotedblright{}.) Say it once more.
\end{tcolorbox}
